\chapter{EU-ret og international ret}

\chapterlead{Dansk ret eksisterer i et lagdelt retssystem. EU-retten kan have direkte betydning i Danmark, mens internationale konventioner får virkning gennem traktat, lovgivning og fortolkning. Den juridiske opgave er at finde det rigtige lag og forstå forbindelsen mellem dem.}

\section{Hvorfor er EU-ret relevant?}

EU-retten regulerer områder som det indre marked, forbrugerbeskyttelse, databeskyttelse, konkurrence, miljø, migration og digital regulering. En dansk sag kan derfor være underlagt både en dansk lov og en EU-forordning eller et direktiv.

EU-rettens vigtigste retskilder er traktaterne, forordninger, direktiver, afgørelser, Domstolen for Den Europæiske Unions praksis og visse generelle principper. EU-charteret beskytter grundlæggende rettigheder, når medlemsstaterne gennemfører eller anvender EU-ret.

\section{Forordninger og direktiver}

En EU-forordning gælder som udgangspunkt direkte i medlemsstaterne. Den skal ikke først omskrives til en dansk lov for at have virkning, selv om national lovgivning kan være nødvendig for at organisere myndigheder, sanktioner og procedurer.

Et direktiv fastlægger et resultat, som medlemsstaterne skal nå, men overlader ofte valg af form og midler til national gennemførelse. Man skal derfor finde både direktivet og den danske gennemførelseslov, når man arbejder med en direktivbaseret regel.

\begin{tabularx}{\textwidth}{@{}>{\RaggedRight\arraybackslash}p{2.7cm}X X@{}}
\toprule
\textbf{Instrument} & \textbf{Grundidé} & \textbf{Spørgsmål i en konkret sag} \\
\midrule
Forordning & Ensartede regler, der som udgangspunkt gælder direkte. & Hvilke bestemmelser og overgangsregler gælder, og hvilken national myndighed håndhæver dem? \\
Direktiv & Bindende mål, der normalt kræver national gennemførelse. & Hvilken dansk regel gennemfører direktivet, og skal den fortolkes direktivkonformt? \\
Traktat & EU's overordnede kompetencer, institutioner og principper. & Er området omfattet af EU's kompetence, og hvilken traktatbestemmelse er relevant? \\
Domstols-\\praksis & Fortolker og konkretiserer EU-retten. & Hvilken retlig test og hvilke betingelser har EU-Domstolen fastlagt? \\
\bottomrule
\end{tabularx}

\section{Forrang og effektivitet}

EU-rettens forrang betyder i hovedtræk, at modstridende national ret må vige inden for EU-rettens anvendelsesområde. Forrang er ikke det samme som, at EU kan regulere alle emner. Først skal man finde en gyldig EU-kompetence og den relevante EU-regel.

National ret skal så vidt muligt fortolkes i overensstemmelse med EU-retten. Hvis en konflikt ikke kan løses gennem fortolkning, kan den nationale domstol skulle undlade at anvende den modstridende nationale regel i den konkrete sag. Den præcise virkning afhænger af regeltype, parternes relation og EU-rettens betingelser.

\section{Præjudicielle spørgsmål}

Hvis en dansk domstol er i tvivl om fortolkningen eller gyldigheden af EU-retten, kan den forelægge spørgsmålet for EU-Domstolen. En forelæggelse er ikke en almindelig appel. Den nationale ret beholder ansvaret for at afgøre den konkrete sag, efter at EU-Domstolen har besvaret det EU-retlige spørgsmål.

Metodisk er pointen: Hvis resultatet afhænger af et uklart EU-retligt begreb, skal man søge EU-Domstolens praksis og ikke kun læse den danske lov.

\section{International ret og dansk ret}

Internationale traktater binder Danmark på folkeretligt plan. Den nationale virkning afhænger af, hvordan traktaten er gennemført eller indarbejdet i dansk ret. Nogle konventioner får betydning gennem en dansk gennemførelseslov; andre påvirker fortolkningen af national ret eller statens internationale ansvar.

Den Europæiske Menneskerettighedskonvention er inkorporeret i dansk ret og har stor betydning for domstole og myndigheder \citep{echr,menneskerettighedsloven}. Den Europæiske Menneskerettighedsdomstols praksis bruges til at forstå konventionens dynamiske anvendelse, men enhver sag kræver opmærksomhed på den konkrete artikel, den konkrete indgrebsform og statens skønsmargin.

\section{Et EU-retligt arbejdsskema}

\practice{Når en sag kan have EU-retlige elementer, så spørg: (1) Er faktum inden for et EU-reguleret område? (2) Hvilken traktatbestemmelse eller EU-retsakt er relevant? (3) Er det en forordning, et direktiv eller en anden retskilde? (4) Findes der dansk gennemførelseslov? (5) Hvad siger EU-Domstolens praksis? (6) Er Charteret relevant? (7) Er der en konflikt med national ret?}

\section{Eksempel: persondata}

En dansk virksomhed bruger en amerikansk cloudtjeneste til at behandle kundedata. Spørgsmålet er ikke kun, om virksomheden følger dansk databeskyttelseslov. Man må også anvende GDPR, vurdere rollerne som dataansvarlig og databehandler, undersøge overførsel til tredjelande og se på eventuelle supplerende danske regler. Hvis teknologien bruger kunstig intelligens, kan AI-forordningen give yderligere pligter, afhængigt af systemets anvendelse, aktørens rolle og den relevante anvendelsesdato.

\section{EU-rettens begrænsninger}

Det er en fejl at bruge EU-retten som et generelt argument for alt. EU har kun de kompetencer, som medlemsstaterne har overladt gennem traktaterne. En sag uden tilstrækkelig forbindelse til EU-rettens område kan ikke automatisk bringes ind under Charteret eller EU's regler.

Det er også en fejl at antage, at alle EU-regler er ensartede. Medlemsstaterne kan have forskellige institutioner, sanktioner og procedurer, så længe de overholder EU-rettens krav om blandt andet effektivitet, ækvivalens og loyal gennemførelse.

\question{En dansk kommune bruger en EU-forordning som begrundelse for et indgreb, men kan ikke pege på den konkrete artikel. Hvad skal du undersøge, før du accepterer henvisningen?}

\section*{Kapitelopsamling}
\begin{itemize}
  \item EU-retten er et integreret lag i mange danske retsområder.
  \item Forordninger gælder som udgangspunkt direkte; direktiver kræver normalt national gennemførelse.
  \item EU-rettens forrang gælder inden for EU-rettens anvendelsesområde.
  \item Internationale konventioners nationale virkning afhænger af traktatens status og dansk gennemførelse.
  \item Begynd altid med kompetence og den konkrete retskilde, ikke med en abstrakt henvisning til ``EU''.
\end{itemize}
